\section{Conclusion}
\label{sec:conclusion}

We introduced a probe-free method for extracting and characterizing the implicit vocabulary of large language models.
By combining representation shift, compositional deviation, and internal similarity signals from model hidden states, our method generalizes to any transformer without requiring pre-trained probes.
Applied to \mistral on 150 Wikipedia articles, we extracted 6{,}494 unique multi-token implicit vocabulary items and showed that items with higher erasure-like scores are significantly less compositional (Spearman $\rho = -0.182$, $p = 0.003$; Cohen's $d = -0.60$).
Named entities constitute the largest non-compositional category (38\% of meaningful items), followed by technical terms (15\%) and fixed expressions (12\%), and the compositionality distribution is continuous rather than binary.

These findings demonstrate that LLMs develop vocabulary-like internal representations that extend far beyond their explicit tokenizer vocabulary, particularly for named entities and conventionalized phrases.
Future work should validate these findings across model families and vocabulary sizes, replace LLM-based compositionality judgments with human annotations, and explore whether the implicit vocabulary can directly inform tokenizer design for improved language modeling.
